

\subsection{模型特点}
本文针对赛题给定的信息时序约束、风险费用结构与结算规则，完成如下建模处理：
\begin{enumerate}
    \item 四个子问题采用统一物理约束与决策可用信息集$\mathcal{I}_{d,t}$，仅根据题目要求修改购电量调整权限与电价时序；实测真值在当前时段决策生成之后，方可用于功率修正与事后费用结算。
    \item 5倍应急购电溢价使得风险裕量具备经济必要性；计划购电与应急购电之间的最优权衡由优化命题给出，理论分位取0.8，实际依托历史净负荷误差的经验分位数实现，并基于该参数完成全年逐日仿真。采用循环能量恒等式独立校验仿真计算结果，分别记录每个时段的0:00原始计划购电量、日内最终调整购电量、弃购电量与紧急购电量，各时段最终购电方案仅相对0:00原始计划执行一次结算。
    \item 设计八组预报信息组合开展费用对比，借助边际收益与$\kappa^*$衡量各预报时刻对应的经济价值；设置理想电价信息基准场景，用于量化电价未知带来的额外购电成本。
    \item 问题三、四采用同一套滚动调度框架，对比不同预报方案与合同调整策略；通过预报组合对照，定量刻画不同预报时刻下总费用、应急购电量的变化规律。
\end{enumerate}

\subsection{局限与推广}
题目未给出并网容量上限以及售电相关机制，本文假设外购电无上限，且微网不可向主网反向送电。若工程场景存在并网容量约束，则需增加约束$g_t+u_t \le G_{\text{max}}$。模型采用经验分位数刻画历史残差尾部特征，未刻画连续多时段预测误差间的相关性，因此逐时段覆盖率不能等价为整日可靠概率。仿真采用10分钟聚合时间粒度，忽略时隙内部测量延迟与功率爬坡过程；实际工程应用时，可缩短仿真步长或增加延迟状态变量予以修正。本文全部结论基于2025年历史数据得到，更换其他年份数据集时，需要重新校验模型参数与费用表现。

按预报提前时间分组设置风险裕量时，采用同一种风险裕量的全局方案，在样本上的总费用相比分组方案低8860.56元。本文仍然选用分组方案，原因在于该方案保留了“预报提前期越长，预测误差通常越大”的数据固有特征。问题三中$H=1$对应的样本费用略低于$H=6$，说明滚动时域长度需要在样本经济性与多时段规划能力之间做折中取舍。

模型可扩展到含风电、可控柔性负荷与多组储能的园区微网场景。新增分布式风电资源，只需在供需平衡方程内增加对应出力供给变量；引入需求响应资源，则增加负荷平移变量，并配套用户用电舒适度约束。若放开售电权限或存在负电价场景，应当采用带有显式充放电互斥约束的混合整数规划模型，规避充放电循环套利问题。微网分层控制的一般性理论背景可参见参考文献。

若能够获取多年实测数据，可将当前35日滑动分位替换为季节条件分位或场景树方法，并采用独立年份样本完成参数$W,\alpha,H$的优选；如果预报供应商按次收取预报费用，可将$\kappa$纳入预报获取方案优化，把盈亏平衡分析拓展为0-1整数选择模型；当园区配备多块储能电池时，可以对每一块电池单独建立SOC状态方程，额外增加共享并网容量约束与电池寿命折损成本。上述所有拓展场景，均可沿用本文的购电计划生成、日内实际调度以及费用结算完整流程。
